\documentclass[11pt,a4paper]{article}
\usepackage[utf8]{inputenc}
\usepackage[T1]{fontenc}
\usepackage{amsfonts}
\usepackage[french]{babel}
\frenchbsetup{StandardLists=true} % à inclure si on utilise \usepackage[francais]{babel}
\usepackage{enumitem}
\usepackage{amssymb}
\usepackage{amsmath}
\usepackage[left=2cm,right=2cm,top=2cm,bottom=2cm]{geometry}
\usepackage{multirow}
\usepackage[dvipsnames]{xcolor}
\usepackage{gensymb} %symbole degré

\usepackage[T1]{fontenc}
\usepackage{fouriernc}
\usepackage{graphicx}

\usepackage{setspace}
\singlespacing

\usepackage{pstricks}
\usepackage{xkeyval}
\usepackage{pst-labo}
\usepackage{chemist} %chimie

\usepackage{tcolorbox}
\usepackage{multicol}
\usepackage{caption}
\usepackage{fancyhdr}
\pagestyle{fancy}
\renewcommand\headrulewidth{1pt}
\fancyhead[L]{Lycée Denis Diderot }
\fancyhead[R]{Classe de 1 Bac Pro }
\setcounter{page}{16}\fancyfoot[C]{page \thepage}
\fancyfoot[R]{Année 2024-2025}
\fancyfoot[L]{Alexis RUIZ}
\usepackage{multirow,tabularx,booktabs}
\usepackage[tikz]{bclogo} 

\newcommand{\cadreblanc}[1]{%
\medskip\framebox[\linewidth][c]{%
\rule{0mm}{#1mm}}\par
\medskip}

\usepackage{awesomebox}
\setlength{\aweboxvskip}{0mm}
\usepackage{pifont}
\usepackage{chemmacros}
\chemsetup{modules=all}
\setlength{\columnseprule}{.25mm}

\renewcommand{\arraystretch}{2.5}
\newcommand*\acocher{\quitvmode{\fboxsep0pt \fboxrule0.8pt \fbox{\vrule height1.8ex depth.1ex width0pt \vrule height0pt depth0pt width1.9ex }}}

\newcolumntype{C}[1]{>{\centering\arraybackslash }m{#1}}

\newcommand{\code}{\awesomebox[black]{2pt}{

\includegraphics[scale=0.15]{Chapitre 1 - Energie et puissance électrique/images/frame.png}
}{black}}

\newcommand{\moteur}{\awesomebox[black]{2pt}{

\includegraphics[scale=1]{images/moteur.jpg} 
}{black}}

\newcommand{\defconclusion}{\awesomebox[black]{2pt}{\rotatebox{90}{\Large{\textbf{Conclusion}}}}{black}}

\newcommand{\defbox}{\awesomebox[black]{2pt}{
\includegraphics[scale=0.5]{Chapitre 1 - Energie et puissance électrique/images/appel exam.png}}{black}}

  \newenvironment{itemize*}%  % listes compactes
    {\begin{termite}%
      \setlength{\itemsep}{0pt}%
      \setlength{\parskip}{0pt}}%
    {\end{itemize}}

\frenchbsetup{StandardLists=true}



\begin{document}

\noindent

\begin{center}
\LARGE{Évaluation- Solutions aqueuses}\\[8mm]
\end{center}
\normalsize
\hrule 


\flushleft \textbf{NOM : \hfill Classe : \hspace{3cm} ~~\\[2mm]
Prénom:\\[0,4cm]
}
\hrule\
\\[0,2cm]

\vspace {8 mm}


\fbox{\textbf{Exercice 1 : Choisir la (ou les) bonne(s) réponse(s)}}


    \begin{tabular}{p{2mm} p{5.5cm}|C{3cm}|C{3cm}|C{3cm}}
    

        \ding{202} & La quantité de matière s'exprime en  :& $\Box $ ~~mole (mol) & $\Box $ ~~gramme (g)   & $\Box $ ~~kilogramme (kg) \\
         \hline
      \ding{203} & La masse molaire du glucose M(\ch{C_6 H_12 O_6} ) =   & $\Box $ ~~ 180 g/mol  & $\Box $ ~~ 90 g/mol   & $\Box $ ~~ 360 g/mol  \\
        \hline
         \ding{204} & Avec la formule $n=\dfrac{m}{M}$, on obtient m avec  &$\Box $ ~~ $m=n \times M$  &$\Box $ ~~ $m=\dfrac{n}{M}$  & $\Box $ ~~$m=\dfrac{M}{n}$   \\
         \hline
         \ding{205} & Si on dilue 5 fois une solution aqueuse d'hydroxyde de sodium de concentration C = 1 mol/L on obtient une concentration  &$\Box $ ~~  C = 5 mol/L  &$\Box $ ~~  C = 0,1 mol/L & $\Box $ ~~ C = 0,2 mol/L \\

         \hline
         
        \ding{206} & On dilue 5 fois une solution en versant  &$\Box $ ~~ 50 mL dans 100 mL  & $\Box $ ~~ 20 mL dans 200 mL  & $\Box $ ~~ 5 mL dans 25~mL \\

 
        
         \end{tabular}

        
\vspace{3mm}


\fbox{\textbf{Exercice 2  }}

A l'aide de la formule $n=\dfrac{m}{M}$, compléter le tableau suivant : 

\begin{center}

    \renewcommand{\arraystretch}{3}


    \begin{tabular}{|C{0.4cm}C{3cm}|C{4cm}|C{4cm}|}
         \hline
      & $n$ & $m$ & $M$ \\
         \hline 
         \ding{202} &  25 mmol & & 40 g/mol \\
         \hline
         \ding{203} &  & 25 g  & 180 g/mol \\
         \hline 
         \ding{204} & 0,1 mol & 1,2 g  & \\
         \hline
         

    \end{tabular}\end{center}

     avec 1 mmol = 0,001 mol

    \vfill 

\newpage

\fbox{\textbf{Exercice 3 }}

\vspace{-1mm}

\noindent%
\begin{minipage}{.40\textwidth}%
  Choisissez et nommez le matériel  dont vous avez besoin pour préparer une disolution. 

\end{minipage}%
\hfill
\begin{minipage}{.55\textwidth}%
\includegraphics[scale=0.7]{Chapitre 4 - Solutions aqueuses/images/materiel.png}
\end{minipage}%

\cadreblanc{25}

\vfill 


\fbox{\textbf{Exercice 4}}
\begin{center}
    
\textbf{Le lait est-il bien conservé ? }
\end{center}
\noindent%

\begin{minipage}{.70\textwidth}%
Léna, responsable qualité, est chargée d'effectuer le titrage de l'acide lactique (\ch{C3H6O3}) contenu dans du lait, afin d’en apprécier l'état de conservation. Elle prélève un volume \ch{V_a} =20 mL de lait qu'elle dilue en ajoutant 180 mL d’eau distillée et quelques gouttes de BBT (indicateur coloré). Elle dose cette solution par une solution d’hydroxyde de sodium de concentration \ch{C_b} = 0,05 mol/L. Après avoir versé un volume équivalent \ch{V_E} =11,3 mL d’hydroxyde de sodium, la solution se colore en bleu de façon persistante, le volume équivalent est donc atteint. .

\end{minipage}%
\hfill
\begin{minipage}{.25\textwidth}%
\fbox{\includegraphics[width=\textwidth]{Chapitre 4 - Solutions aqueuses/lait.png}}
\end{minipage}%

\vspace{5mm}

\textbf{\ding{202} Réaliser} \\

En utilisant la relation $C_a \times V_a = C_b \times V_E$ , calculer, en mol/L, la concentration en acide lactique contenue dans le lait. \\

\cadreblanc{15}











\textbf{\ding{203} Valider} \\

 Le lait ne doit pas contenir plus de 0,024 mol/L d'acide lactique. Le lait dosé par Léna a-t-il été convenablement conservé?

\cadreblanc{15}



\textbf{\ding{204} Communiquer} \\[3mm]

La dilution du lait permet de mieux apprécier le virage de l’indicateur. Expliquez pourquoi elle ne modifie pas la valeur du volume équivalent \ch{V_E}.

\cadreblanc{15}



  \vfill 
  

\end{document}